% Part: first-order-logic
% Chapter: completeness
% Section: compactness-direct

\documentclass[../../../include/open-logic-section]{subfiles}

\begin{document}

\iftag{FOL}
      {\olfileid[hi]{fol}{com}{cpd}}
      {\olfileid[hi]{pl}{com}{cpd}}

\olsection{संहतता प्रमेय का प्रत्यक्ष प्रमाण}

पूर्णता प्रमेय का सहारा लिए बिना, पूर्णता प्रमेय के प्रमाण वाले ही विचारों
से हम संहतता प्रमेय को सीधे सिद्ध कर सकते हैं। पूर्णता प्रमेय के प्रमाण में
हमने संगत समुच्चय~$\Gamma$ से आरंभ किया, जो !!{sentence}s का था; उसे
संगत\iftag{FOL}{, संतृप्त}{} और !!{complete} समुच्चय~$\Gamma^*$ तक
विस्तारित किया, जो !!{sentence}s का था; और फिर दिखाया कि
\iftag{FOL}{पद-प्रतिरूप~$\Struct{M(\Gamma^*)}$}{!!{valuation}~$\pAssign{v(\Gamma^*)}$}
में, जिसे~$\Gamma^*$ से बनाया गया था,~$\Gamma$ के सभी !!{sentence}s सत्य
हैं; अतः $\Gamma$ संतुष्टनीय है।

इसी विधि से हम दिखा सकते हैं कि परिमिततः संतुष्टनीय वाक्य-समुच्चय
संतुष्टनीय है। हमें सत्यता उपपत्ति तक ले जाने वाले परिणामों के अनुरूप रूपों
को बस इस प्रकार सिद्ध करना है कि उनमें ``संगत'' के स्थान पर ``परिमिततः
संतुष्टनीय'' हो।

\begin{prop}
\ollabel{prop:fsat-ccs}
मान लें कि $\Gamma$ !!{complete} और परिमिततः संतुष्टनीय है। तब:
\begin{enumerate}
\tagitem{prvAnd}{$(!A \land !B) \in \Gamma$
  तभी और केवल तभी जब $!A \in \Gamma$ और $!B \in \Gamma$ दोनों हों।}{}

\tagitem{prvOr}{$(!A \lor !B) \in \Gamma$ तभी और केवल तभी जब
  या तो $!A \in \Gamma$ या $!B \in \Gamma$ हो।}{}

\tagitem{prvIf}{$(!A \lif !B) \in \Gamma$ तभी और केवल तभी जब
  या तो $!A \notin \Gamma$ या $!B \in \Gamma$ हो।}{}
\end{enumerate}
\end{prop}

\tagprob{FOL}
\begin{prob}
\olref[fol][com][cpd]{prop:fsat-ccs} सिद्ध कीजिए। $\Proves$ का उपयोग न करें।
\end{prob}
\tagendprob

\tagprob{notFOL}
\begin{prob}
\olref[pl][com][cpd]{prop:fsat-ccs} सिद्ध कीजिए। $\Proves$ का उपयोग न करें।
\end{prob}
\tagendprob

\iftag{FOL}{%
\begin{lem}
\ollabel{lem:fsat-henkin} प्रत्येक परिमिततः संतुष्टनीय समुच्चय~$\Gamma$ को
किसी संतृप्त परिमिततः संतुष्टनीय समुच्चय~$\Gamma'$ तक विस्तारित किया जा
सकता है।
\end{lem}
}{}

\tagprob{FOL}
\begin{prob}
\olref[fol][com][cpd]{lem:fsat-henkin} सिद्ध कीजिए। (संकेत: निर्णायक चरण यह
दिखाना है कि यदि $\Gamma_n$ परिमिततः संतुष्टनीय है, तो $\Gamma_n
\cup \{!D_n\}$ भी परिमिततः संतुष्टनीय है, और इसमें !!{derivation}s अथवा
संगतता का कोई सहारा न लिया जाए।)
\end{prob}
\tagendprob

\iftag{FOL}{
\begin{prop}\ollabel{prop:fsat-instances}
मान लें कि $\Gamma$ पूर्ण, परिमिततः संतुष्टनीय और संतृप्त है।
\begin{tagenumerate}{prvEx,prvAll}
\tagitem{prvEx}{$\lexists[x][!A(x)] \in \Gamma$ तभी और केवल तभी जब
  $!A(t) \in \Gamma$ कम-से-कम एक बंद पद~$t$ के लिए हो।}{}
\tagitem{prvAll}{$\lforall[x][!A(x)] \in \Gamma$ तभी और केवल तभी जब
  $!A(t) \in \Gamma$ सभी बंद पदों~$t$ के लिए हो।}{}
\end{tagenumerate}
\end{prop}
}{}

\tagprob{FOL}
\begin{prob}
\olref[fol][com][cpd]{prop:fsat-instances} सिद्ध कीजिए।
\end{prob}
\tagendprob

\begin{lem}
\ollabel{lem:fsat-lindenbaum} प्रत्येक परिमिततः संतुष्टनीय समुच्चय~$\Gamma$
को !!a{complete} और परिमिततः संतुष्टनीय समुच्चय~$\Gamma^*$ तक विस्तारित
किया जा सकता है।
\end{lem}

\tagprob{FOL}
\begin{prob}
\olref[fol][com][cpd]{lem:fsat-lindenbaum} सिद्ध कीजिए। (संकेत: निर्णायक
चरण यह दिखाना है कि यदि $\Gamma_n$ परिमिततः संतुष्टनीय है, तो
$\Gamma_n \cup \{!A_n\}$ या $\Gamma_n \cup \{\lnot !A_n\}$ में से कोई एक
परिमिततः संतुष्टनीय है।)
\end{prob}
\tagendprob

\tagprob{notFOL}
\begin{prob}
\olref[pl][com][cpd]{lem:fsat-lindenbaum} सिद्ध कीजिए। (संकेत: निर्णायक
चरण यह दिखाना है कि यदि $\Gamma_n$ परिमिततः संतुष्टनीय है, तो
$\Gamma_n \cup \{!A_n\}$ या $\Gamma_n \cup \{\lnot !A_n\}$ में से कोई एक
परिमिततः संतुष्टनीय है।)
\end{prob}
\tagendprob

\begin{thm}[संहतता]
\ollabel{thm:compactness-direct} $\Gamma$ संतुष्टनीय तभी और केवल तभी है जब
वह परिमिततः संतुष्टनीय है।
\end{thm}

\begin{proof}
यदि $\Gamma$ संतुष्टनीय है, तो ऐसी
\iftag{FOL}{!!a{structure}~$\Struct{M}$}{!!a{valuation}~$\pAssign{v}$} है कि
$\iftag{FOL}{\Sat{M}}{\pSat{v}}{!A}$ सभी $!A \in \Gamma$ के लिए। निस्संदेह,
यह \iftag{FOL}{$\Struct{M}$}{$\pAssign{v}$} भी~$\Gamma$ के प्रत्येक परिमित
उपसमुच्चय को संतुष्ट करती है, अतः $\Gamma$ परिमिततः संतुष्टनीय है।

अब मान लें कि $\Gamma$ परिमिततः संतुष्टनीय है।\iftag{FOL}{
  \olref{lem:fsat-henkin} से कोई परिमिततः संतुष्टनीय, संतृप्त समुच्चय
  $\Gamma' \supseteq \Gamma$ है।}{} \olref{lem:fsat-lindenbaum} से
$\iftag{FOL}{\Gamma'}{\Gamma}$ को !!a{complete} और परिमिततः संतुष्टनीय
समुच्चय~$\Gamma^*$ तक विस्तारित किया जा सकता है\iftag{FOL}{, और
  $\Gamma^*$ अब भी संतृप्त है}{}। \olref[mod]{defn:termmodel} की तरह
\iftag{FOL}{पद-प्रतिरूप~$\Struct{M(\Gamma^*)}$}{!!{valuation}~$\pAssign{v(\Gamma^*)}$}
बनाएँ।\iftag{FOL}{ ध्यान दें कि \olref[mod]{prop:quant-termmodel} ने इस तथ्य
पर भरोसा नहीं किया था कि $\Gamma^*$ संगत है (या !!{complete} अथवा संतृप्त,
इस विषय में), बल्कि केवल इस तथ्य पर कि $\Struct{M(\Gamma^*)}$ आच्छादित है।}{}
सत्यता उपपत्ति (\olref[mod]{lem:truth}) का प्रमाण तब भी चलता है जब हम
\olref[ccs]{prop:ccs}\iftag{FOL}{ और
\olref[hen]{prop:saturated-instances} के संदर्भों के स्थान पर क्रमशः
\olref{prop:fsat-ccs} और \olref{prop:fsat-instances} के संदर्भ रखें}।
\end{proof}

\tagprob{FOL}
\begin{prob}
सत्यता उपपत्ति \olref[fol][com][mod]{lem:truth} का वह पूरा प्रमाण लिखिए जो
\olref[fol][com][cpd]{thm:compactness-direct} के प्रमाण के लिए आवश्यक है।
\end{prob}
\tagendprob

\tagprob{notFOL}
\begin{prob}
सत्यता उपपत्ति \olref[pl][com][mod]{lem:truth} का वह पूरा प्रमाण लिखिए जो
\olref[pl][com][cpd]{thm:compactness-direct} के प्रमाण के लिए आवश्यक है।
\end{prob}
\tagendprob

\end{document}
